% Three routes to a torque figure, and which of them this bench can run.
\begin{tikzpicture}[font=\sffamily\small, x=1cm, y=1cm]

% ---------------- route 1: a sensor in the joint
\node[chlabel, anchor=west] at (-6.6,4.4) {route 1: measure it};
\node[absentblk, text width=36mm] (sens) at (-4.6,3.4)
  {a sensing element in the joint\\{\scriptsize strain gauges on a flexure, measuring the transmitted torque}};
\node[absentblk, text width=36mm] (cond) at (-4.6,1.4)
  {its conditioning chain\\{\scriptsize excitation, a chopper amplifier, a 24-bit converter, a calibration}};
\node[chlabel, anchor=north, text width=38mm, align=center] at (-4.6,0.3)
  {the honest answer, and the expensive one};

% ---------------- route 2: current
\node[chlabel, anchor=west] at (-1.4,4.4) {route 2: infer it from current};
\node[absentblk, text width=36mm] (cur) at (0.6,3.4)
  {current sensing\\{\scriptsize a shunt, an amplifier, and a converter synchronised to the switching}};
\node[absentblk, text width=36mm] (gear) at (0.6,1.4)
  {and a gearbox model\\{\scriptsize because what the motor makes is not what the joint delivers}};
\node[chlabel, anchor=north, text width=38mm, align=center] at (0.6,0.3)
  {what most drives already have, and wrong by the losses};

% ---------------- route 3: the observer
\node[chlabel, anchor=west] at (3.6,4.4) {route 3: infer it from motion};
\node[fw, text width=36mm] (pos) at (5.8,3.4)
  {position, from the real decode path\\{\scriptsize chapter 5, and the only real input here}};
\node[modelled, text width=36mm] (obs) at (5.8,1.4)
  {the observer\\{\scriptsize J times acceleration, minus the commanded torque, plus the friction model}};
\node[chlabel, anchor=north, text width=38mm, align=center] at (5.8,0.3)
  {what this chapter builds, and what industrial machines use with no sensor};

% ---------------- what they all feed
\node[user, text width=34mm] (est) at (0.6,-1.6)
  {the estimate, with its provenance\\{\scriptsize value, source, model version, filter latency}};
\node[user, text width=34mm] (det) at (0.6,-3.6)
  {the contact detector\\{\scriptsize threshold, hold, latch, and a measured false-alarm rate}};

\draw[flow] (sens) -- (cond);
\draw[flow] (cur) -- (gear);
\draw[flow] (pos) -- (obs);
\draw[hiz, ->] (cond.south) -- ++(0,-0.7) -| (est.west);
\draw[hiz, ->] (gear) -- (est);
\draw[flow] (obs.south) -- ++(0,-0.7) -| (est.east);
\draw[flow] (est) -- (det);

% ---------------- the one real measurement in the chapter
\node[hw, text width=34mm] (imu) at (6.4,-3.6)
  {the inertial unit\\{\scriptsize chapter 6. Real silicon, and a tap on the bench is a real impulse}};
\draw[flow] (imu) -- node[lbl, above]{a different question} (det);

\node[umlnote, text width=52mm, anchor=north west] at (-7.4,-1.4)
  {Two of the three routes need hardware this bench does not have, and the
   figure says what each would require rather than leaving the reader to
   imagine it. The third needs no sensor at all, which is why it is what
   industrial machines fall back on, and why its error budget is the whole of
   this chapter.};
\end{tikzpicture}
